\subsection{The actual relative exponential with its prescribed central fibre}
\label{subsec:ftb-actual-relative-exponential}

Here \emph{holomorphic} for a map into scheme points means
\(\mathrm{IsHolTo}\): the inverse image of each algebraic open
is open, and every regular-function evaluation on that
inverse image is holomorphic.
All points and all group operations below are those of the
given algebraic family.

\begin{srlemma}[Relative affine coordinates with finite inverse tests]
\label{lem:ftb-relative-affine-coordinates}
Let \(p:Z\to X\) be smooth of relative dimension \(g\).
Let \(B\subset\mathbb C^m\) be open, let
\(y:B\to X(\mathbb C)\) be regular-evaluation-holomorphic,
and let \(z_0\in Z(\mathbb C)\) lie over \(y(b_0)\).
There are an affine neighborhood \(U\) of \(z_0\), regular
functions \(t_1,\ldots,t_g\) on \(U\), finitely many further
regular functions \(f_1,\ldots,f_n\), open parameter-coordinate
sets \(N\subset\mathbb C^m\times\mathbb C^g\),
\(O\subset\mathbb C^m\times\mathbb C^n\), and a
regular-evaluation-holomorphic \(\sigma:N\to U(\mathbb C)\)
such that
\[
 p(\sigma(b,w))=y(b),\qquad t(\sigma(b,w))=w.                 \tag{RE1}
\]
The point \((b_0,f(z_0))\) belongs to \(O\).
If \(a\in U(\mathbb C)\) lies over \(y(b)\) and
\((b,f(a))\in O\), then
\[
                (b,t(a))\in N,\qquad \sigma(b,t(a))=a.       \tag{RE2}
\]
Moreover, two dual-number points in \(U\), with the same
base point, the same reduced point and the same values of
all \(t_i\), are equal.
\end{srlemma}
\begin{proof}
Take affine neighborhoods \(V\subset X\), \(U\subset Z\)
on which smoothness gives an etale ring homomorphism
\[
                A[T_1,\ldots,T_g]\longrightarrow R,
       \quad A=\Gamma(V,\mathcal O_X),\
             R=\Gamma(U,\mathcal O_Z).
\]
The \(t_i\) are the images of the \(T_i\).
Shrink the parameter neighborhood so that \(y(b)\in V\);
evaluation along \(y\) gives ring maps
\(\varphi_b:A\to\mathbb C\), analytic on every coefficient.

Use a standard smooth presentation of this etale algebra.
There are finitely many auxiliary variables \(Y_j\),
as many equations \(P_r(T,Y)\), and an invertible square
\(Y\)-Jacobian at the given evaluation point.
The map
\[
 (b,w,v)\longmapsto
 \bigl(b,w,(P_r^{\varphi_b}(w,v))_r\bigr)                    \tag{RE3}
\]
is analytic.
Its derivative is block triangular, with identity on
\((b,w)\) and that invertible Jacobian on \(v\).
The analytic inverse theorem gives an analytic inverse
on actual open neighborhoods of the point and its image.
At target equation-coordinate \(0\), its \(Y\)-coordinates
solve the relations.
Evaluation therefore factors through the quotient
presentation and gives a ring homomorphism
\(\sigma(b,w):R\to\mathbb C\).

For each \(r\in R\), choose a polynomial representative
in the presentation.
Its evaluation at these analytic coordinates is analytic.
Thus every regular-function value is analytic.
On smaller algebraic opens a regular function is locally
a quotient of such values with nonzero denominator;
this proves the claimed \(\mathrm{IsHolTo}\) statement.

For the inverse tests choose the finite list consisting
of the \(t_i\) and the auxiliary generators \(Y_j\).
Let \(O\) say that their values, together with \(b\),
lie in the source of the inverse in (RE3).
Every ring homomorphism in this source that satisfies
the equations is the indicated inverse point, by the
two-sided inverse identities in (RE3).
This gives (RE2), including actual equality of the
ring homomorphisms, hence of the scheme points.
Finally, formal unramifiedness of the etale ring map
says that two lifts across
\(\mathbb C[\epsilon]\to\mathbb C\), agreeing on \(A[T]\)
and after reduction, coincide.
The hypotheses on the base and on the \(t_i\) are exactly
this agreement.
\end{proof}

\begin{srlemma}[The parameter-dependent halving logarithm]
\label{lem:ftb-parametric-halving-logarithm}
Let \(M(b,x,y)\), with \(b\in\mathbb C^m\) and
\(x,y\in\mathbb C^g\), be analytic near \((b_0,0,0)\).
Suppose the unit, commutativity and associativity identities
\[
 M(b,x,0)=x,\quad M(b,x,y)=M(b,y,x),\quad
 M(b,M(b,x,y),z)=M(b,x,M(b,y,z))
                                                               \tag{RE4}
\]
hold as germs at the indicated zero tuples.
There are jointly analytic germs \(E(b,v),L(b,x)\) near
\((b_0,0)\), inverse in their second variables, with
\(E(b,0)=L(b,0)=0\), \(d_vE(b,0)=d_xL(b,0)=I\),
and
\[
       E(b,v+w)=M(b,E(b,v),E(b,w))                            \tag{RE5}
\]
as a joint germ.
\end{srlemma}
\begin{proof}
The inverse theorem with parameters follows by applying
the analytic inverse theorem to \((b,x)\mapsto(b,F(b,x))\).
If \(d_xF(b_0,0)\) is injective, finite equal dimensions
make it invertible; the full derivative is block triangular
and invertible.
The inverse's first coordinate is \(b\), by its right
inverse identity.
This gives a jointly analytic inverse in the second variable.

Apply this to doubling \(x\mapsto M(b,x,x)\).
The derivative in \(x\) is \(2I\), by the two unit identities,
so its inverse \(\varepsilon(b,x)\) is jointly analytic.
For each nearby \(b\),
\[
                 \varepsilon_b(x)=\tfrac12x+O(\|x\|^2).
                                                               \tag{RE6}
\]
The source makes this estimate uniform as follows.
Choose a common product ball of radius \(s>0\) on which
\(\varepsilon\) is analytic, \(\|\varepsilon\|\le K\),
and its slices have (RE6), with \(K>0\).
The second-order Cauchy estimate, with
\(\varepsilon_b(0)=0\), \(d\varepsilon_b(0)=I/2\),
gives
\[
 \|\varepsilon_b(x)-x/2\|\le C\|x\|^2,\qquad C=4K/s^2.
\]
Put \(r=\min(s,1/(8C))\).
For \(\|x\|<r\) this implies
\[
                       \|\varepsilon_b(x)\|\le(5/8)\|x\|.
                                                               \tag{RE7}
\]
All iterates stay in that ball.
The functions
\(L_j(b,x)=2^j\varepsilon_b^{\,j}(x)\) are jointly analytic.
The contraction and quadratic-error calculation used in
Lemma~\ref{lem:ftb-halving-logarithm} gives the uniform tail
\[
 \|L_j(b,x)-L(b,x)\|
    \le \frac{64}{7}C\|x\|^2(25/32)^j
    \le \frac{64}{7}Cr^2(25/32)^j.                          \tag{RE8}
\]
Thus \(L\) is jointly analytic by locally uniform
convergence of holomorphic functions.

For each fixed nearby \(b\), the same local commutative
group-law argument as in that lemma gives
\(L_b(0)=0\), \(dL_b(0)=I\), and
\(L_b(M_b(x,y))=L_b(x)+L_b(y)\) as a slice germ.
One must also make the last identity uniform in \(b\).
The difference of its sides is jointly analytic on a
common product ball.
For every nearby \(b\) it vanishes on some neighborhood
of \(0\) in the \((x,y)\)-slice.
The identity theorem on the fixed connected slice ball
makes it vanish on that whole ball.
Hence it vanishes as a joint germ.
Apply the parameter inverse theorem once more to \(L\)
to obtain \(E\).
The inverse identities give \(E_b(0)=0\), \(dE_b(0)=I\);
the just-proved joint logarithm identity gives (RE5).
\end{proof}

\begin{srproposition}[The actual relative exponential]
\label{prop:ftb-actual-relative-exponential}
Let \(k\) be a field with an embedding
\(\iota:k\hookrightarrow\mathbb C\).
Let \(X\to\operatorname{Spec}k\) be locally of finite type,
let \(G/X\) be an abelian scheme smooth of relative
dimension \(g\), and let \(B\subset\mathbb C^m\) be open.
Suppose \(y:B\to X(\mathbb C)\) is
regular-evaluation-holomorphic and each \(y(b)\) lies
over the fixed embedding \(\iota\).
At \(b_0\in B\) prescribe an additive surjective
regular-evaluation-holomorphic
\[
                       e_0:\mathbb C^g\to G_{y(b_0)}(\mathbb C)
\]
whose derivative identifies analytic vectors with
algebraic vertical tangent vectors in both directions.
Then, after shrinking to an open \(b_0\in B'\subset B\),
there are maps \(e_b:\mathbb C^g\to G_{y(b)}(\mathbb C)\)
and a single \(r>0\), with the following properties:
\[
\begin{gathered}
 e_{b_0}=e_0,\quad e_b(v+w)=e_b(v)e_b(w),\quad
                         e_b\text{ is surjective};\\
 (b,v)\longmapsto e_b(v)\text{ is regular-evaluation-holomorphic};\\
 \|v\|<r,\ e_b(v)=1\quad\Longrightarrow\quad v=0
                         \qquad(b\in B').                 \tag{RE9}
\end{gathered}
Each \(e_b\) has both vertical tangent-bijection clauses.

The construction also has the following local lifting
property, with no restriction to the identity section.
If \(W\subset\mathbb C^n\) is open,
\(\beta:W\to B'\) is holomorphic,
\(a_w\in G_{y(\beta(w))}(\mathbb C)\) is
regular-evaluation-holomorphic, and
\(e_{\beta(w_0)}(z_0)=a_{w_0}\), there is an analytic
\(\zeta\) near \(w_0\), with
\[
        \zeta(w_0)=z_0,\qquad
        e_{\beta(w)}(\zeta(w))=a_w\quad\text{near }w_0.       \tag{RE10}
\]
\end{srproposition}
\begin{proof}
Apply Lemma~\ref{lem:ftb-relative-affine-coordinates}
to \(G\to X\) at the identity over \(y(b_0)\).
Set \(c(b)=t(1_b)\), analytic by evaluation along the
algebraic identity section, and write
\[
                 t_b^0(a)=t(a)-c(b),\qquad
                 s_b(x)=\sigma(b,x+c(b)).
\]
On the finite-test neighborhood, (RE1)--(RE2) become
\[
          t_b^0(s_b(x))=x,\quad s_b(0)=1_b,\quad
                         s_b(t_b^0(a))=a.                 \tag{RE11}
\]
All three identities hold on common neighborhoods of
the indicated central points.
Products, inverses and powers of actual holomorphic
families of points are holomorphic:
on affine pieces of the fibre product, its coordinate
ring is generated by the two factor rings;
evaluation preserves their sums and products.
Postcomposition with the algebraic group operation
therefore preserves regular-evaluation holomorphicity.
This also explains why the group law used next is the
one of the given family.

Define
\[
                   M(b,x,y)=t_b^0(s_b(x)s_b(y)).
\]
It is jointly analytic near \((b_0,0,0)\).
Identities (RE11) and the actual group laws imply (RE4).
In the associativity identity, first use (RE11) to
replace \(s_b(M(b,x,y))\) by \(s_b(x)s_b(y)\)
on a small common neighborhood; then use associativity
in \(G\).
Lemma~\ref{lem:ftb-parametric-halving-logarithm} gives \(E,L\).
Consequently
\[
                         v\longmapsto s_b(E(b,v))           \tag{RE12}
\]
is locally additive, jointly holomorphic, and has
coordinate derivative \(I\).

Choose one ball of radius \(r_1>0\) and one neighborhood
of \(b_0\) on which these local identities, the inverse
identities for \(E,L\), and (RE11) hold.
The local additive maps (RE12) extend to additive maps
\(\widehat e_b\) on all of \(\mathbb C^g\).
Explicitly, for sufficiently large \(j\), use
\[
           \widehat e_b(v)=
                  s_b(E(b,2^{-j}v))^{\,2^j}.               \tag{RE13}
\]
The local doubling identity makes (RE13) independent
of sufficiently large \(j\).
The local additive identity, after a common division
of two vectors, proves global additivity.
The extension agrees with (RE12) throughout a smaller
common ball, not merely on an unspecified germ:
divide a point in that ball dyadically until it lies in
the initial equality neighborhood and then double back.

Joint holomorphicity also holds at an arbitrary \((b,v)\).
Choose an integer \(h>0\) so that \(\|v/h\|<r_1\).
Near \((b,v)\), additivity gives
\[
                   \widehat e_{b'}(v')=
                            s_{b'}(E(b',v'/h))^{\,h}.
\]
The right side is jointly regular-evaluation-holomorphic.
Thus the global family has the required holomorphicity.

Shrink the parameter set once more so that \(L_b(0)=0\)
and its inverse identities hold on one common coordinate
neighborhood.
For some uniform \(r_2>0\), if \(\|v\|<r_2\) and
\(\widehat e_b(v)=1\), (RE11)--(RE12) give
\(E(b,v)=0\), hence \(v=L(b,0)=0\).
There is also an actual inverse near the identity:
\[
                 \widehat e_b(L(b,t_b^0(a)))=a,             \tag{RE14}
\]
whenever \(a\) satisfies the original finite open tests
and \((b,t_b^0(a))\) lies in the common logarithm domain.
These are finitely many regular-function tests on each
fixed fibre.
The image of the additive map \(\widehat e_b\) is therefore
a subgroup containing such a neighborhood.
The final assertion of
Proposition~\ref{prop:ftb-integral-complex-points-connected}
makes it the whole fibre.
This proves surjectivity without an unproved assertion
about connected analytic images.

We spell out the tangent step.
For a holomorphic fibre map \(v\mapsto a(v)\), a vector \(w\)
defines on an affine identity chart the homomorphism
\[
             f\longmapsto f(a(0))
                    +\epsilon\,d(f\circ a)_0(w)
                \quad\text{into }\mathbb C[\epsilon].      \tag{RE15}
\]
The product rule makes it a ring homomorphism.
The affine point--ring-homomorphism identification makes
(RE15) an actual dual-number point.
Localization shows it has the asserted derivatives on
every open; equality on an affine chart proves uniqueness.
Since \(a\) remains over a fixed \(y(b)\), its base derivative
is zero, so this point is vertical.
For \(a=\widehat e_b\), the coordinate derivative in \(t\)
is \(I\).
Any vertical tangent vector has a coordinate vector
\(w=(dt_i)\).
The jet (RE15) for that \(w\) has the same base, reduction
and \(t_i\)-values as the given tangent vector.
The final, formally unramified clause of
Lemma~\ref{lem:ftb-relative-affine-coordinates}
makes the two equal.
This proves both tangent-bijection clauses for
\(\widehat e_b\).
More generally the same argument shows that such a
holomorphic map has the two-way tangent correspondence
exactly when its \(g\)-by-\(g\) coordinate derivative
is injective, hence bijective.

The family must now match the \emph{prescribed} central
exponential, rather than an exponential with an arbitrary
central normalization.
Put
\[
             T_0=d\bigl(t\circ e_0\bigr)_0
                         \in\operatorname{End}_{\mathbb C}(\mathbb C^g).
\]
The preceding coordinate characterization, applied to the
supplied tangent correspondence for \(e_0\), makes \(T_0\)
invertible.
There is an analytic local lift \(\eta\) of \(e_0\)
through \(\widehat e_{b_0}\), by (RE14), with \(\eta(0)=0\).
Its derivative is \(T_0\).
The difference
\(\eta(v+w)-\eta(v)-\eta(w)\) tends to zero and lies in
\(\ker\widehat e_{b_0}\).
The uniform small-kernel assertion makes it zero near
\((0,0)\).
A continuous locally additive germ agrees with its derivative:
extend it additively by dyadic division; the extension is
real linear by continuity, and uniqueness of the real
derivative identifies it with \(T_0\).
Thus \(\eta=T_0\) near \(0\), so
\(\widehat e_{b_0}T_0=e_0\) near \(0\), hence everywhere
by additivity.
Define
\[
                              e_b(v)=\widehat e_b(T_0v).
\]
This proves the exact central equality in (RE9).
Surjectivity and both tangent clauses persist because
\(T_0\) is invertible.
Replace \(r_2\) by \(r_2/(\|T_0\|+1)\); the operator-norm
inequality proves the uniform last clause of (RE9).

Finally, let \(\beta,a,w_0,z_0\) be as in (RE10).
The actual holomorphic point
\[
                     h_w=a_w\,e_{\beta(w)}(z_0)^{-1}
\]
equals the identity at \(w_0\), hence satisfies the finite
tests and the logarithm-domain condition near \(w_0\).
Formula (RE14) gives the explicit lift
\[
 \zeta(w)=T_0^{-1}L\bigl(\beta(w),
                     t_{\beta(w)}^0(h_w)\bigr)+z_0.        \tag{RE16}
\]
It is analytic near \(w_0\), equals \(z_0\) there, and
additivity plus (RE14) gives (RE10).
This is the source's full local lifting clause.
\end{proof}

\paragraph{Frozen source ledger and next interface.}
All the following files have prefix
\srcid{src/Development/AbelianVarietyTheory/Proofs/ModuliProofs/},
at \srcid{alpoge-source-7cbf7a0c}.
The square-Jacobian construction (RE3), its finite tests
and quotient factorization are
\srcid{ExistsHolomorphicSections.lean}
(\srcid{a288c7894564d55a8a32b1c04d9bcb42acd6d74b}).
Passage to actual scheme points and formal-unramified
dual-number uniqueness is
\srcid{ExistsChartWithParameters.lean}
(\srcid{868938f5f6aac923cd3bb1064cbbdb99fb0565c7}).
The parameter inverse, uniform Cauchy estimate,
\(5/8\) contraction, \(25/32\) geometric ratio,
\(64/7\) tail, and slice-to-joint identity argument are
\srcid{ExistsAnalyticParamInverse.lean}
(\srcid{fa9a9ead5ef46d2371247240442ac56a161d68da}).
The centered chart, actual local group law, global
additive extension, uniform kernel ball and local inverse
(RE14) are
\srcid{AnalyticAtCoordinateZero.lean}
(\srcid{361cb89abffd48f035e4b979aa93dcdba0af4250}).
Its holomorphic product/pullback rules and additive germ
uniqueness are
\srcid{EventuallyEqOfAdditiveOfHasFDerivAt.lean}
(\srcid{cecff5d92fc66f83c4e05e50aeee9e29a322ca8d}).
The actual dual-number tangent construction and injective
coordinate-derivative criterion are
\srcid{VerticalDerivativeInCoordinates.lean}
(\srcid{c407e0b090e2dd8b374aef3aaa0c1c245e280d29}).
The full assembly, prescribed central equality and
(RE16) are \srcid{ExistsRelativeExpOfParts.lean}
(\srcid{7909f077337921a132a7554acdc1599b9c7f9c6d}),
the theorem \srcid{RelExp.exists_relative_exp}.
The connected-fibre supplier it actually invokes was
proved in
Proposition~\ref{prop:ftb-integral-complex-points-connected}.

Thus the actual \srcid{RelExp.exists_relative_exp}
interface, including its general \(\mathrm{HasLifts}\)
clause and one uniform small-kernel radius, is established.
It does not yet choose a common period lattice, transport
the \(\mathcal O_E\)-action or level markings, or choose a
theta frame in a family.  Those subsequent choices must
use these same actual maps \(e_b\).
